\chapter{Manual de código}
\label{ap:manual-codigo}

De acuerdo con la normativa de la EPSC \cite{reglamento2024,guiaGII2025}, el
manual de código puede consistir en una referencia a un repositorio accesible.
El código fuente de este Trabajo Fin de Grado se encuentra disponible de forma
pública en:

\begin{center}
\url{https://github.com/Jsoriano99/tfg-runner-dashboard}
\end{center}

El repositorio contiene la estructura completa del proyecto, el código
organizado por módulos (ingesta de datos, métricas, análisis, modelos
predictivos y dashboard) y su documentación asociada, de modo que cualquier
persona puede consultar, reproducir y auditar el trabajo.

\section*{Instalación y ejecución}
\label{sec:codigo-instalacion}

\pendiente{incluir las instrucciones exactas para clonar el repositorio, crear
el entorno virtual, instalar las dependencias fijadas en \texttt{requirements.txt}
y ejecutar la aplicación.}

\section*{Estructura del proyecto}
\label{sec:codigo-estructura}

\pendiente{describir el árbol de directorios del repositorio y la
responsabilidad de cada módulo, indicando dónde vive la lógica de cada parte del
sistema.}

\section*{Ejecución de las pruebas}
\label{sec:codigo-pruebas}

\pendiente{explicar cómo ejecutar la batería de pruebas con \texttt{pytest},
dónde se encuentran los \emph{fixtures} anonimizados y qué cubre la suite.}

\section*{Versión y etiqueta de la entrega}
\label{sec:codigo-version}

\pendiente{indicar la versión o etiqueta (\emph{tag}) del repositorio que
corresponde exactamente al estado del código entregado junto con esta memoria.}
